% Source: source/普通高考/2011/2011天津文.pdf, page 1, question 3.
% pdfimages -list reports no embedded bitmap; the program flowchart is vector/text artwork.
% Grid evidence: tmp/review_2011_tianjin_liberal/grid/page_grid100.png,
%   tmp/review_2011_tianjin_liberal/grid/page_grid50.png.
% No-grid source crop: tmp/review_2011_tianjin_liberal/crop/q3_orig_fig_final.png.
% Nodes: 开始; 输入 x; |x|>3?; x=|x−3|; y=2^x; 输出 y; 结束.
% Directed edges: 开始→输入→判定; 判定(是)→绝对值更新→左回环→判定;
% 判定(否)→y=2^x→输出 y→结束. The loop return and output branch are independent.
\begin{tikzpicture}[
  >=Stealth,
  line width=0.55pt,
  line cap=round,
  line join=round,
  every node/.style={font=\normalsize,anchor=center,align=center,
    execute at begin node={\setlength{\parindent}{0pt}\noindent}},
  flowbox/.style={draw,minimum height=.68cm,inner xsep=4pt,inner ysep=2pt,
    anchor=center,align=center},
  terminal/.style={flowbox,rounded corners=.18cm,minimum width=1.35cm,
    minimum height=.72cm},
  io/.style={flowbox,shape=trapezium,trapezium left angle=72,
    trapezium right angle=108,minimum height=.78cm},
  decision/.style={flowbox,shape=diamond,aspect=2.80,inner sep=2pt},
  branchlabel/.style={anchor=center,align=center,inner sep=0pt,
    execute at begin node={\setlength{\parindent}{0pt}}},
  flowarrow/.style={->,line width=0.55pt},
]
  \node[terminal] (start) at (0,9.30) {开始};
  \node[io,minimum width=2.35cm] (input) at (0,8.00) {输入 $x$};
  \coordinate (loopjoin) at (0,7.12);
  \node[decision,minimum width=3.55cm,minimum height=1.00cm] (test) at (0,6.28)
    {$|x|>3?$};
  \node[flowbox,minimum width=3.25cm] (absupdate) at (0,4.78)
    {$x=|x-3|$};
  \node[flowbox,minimum width=2.30cm] (power) at (0,3.10)
    {$y=2^x$};
  \node[io,minimum width=2.05cm] (output) at (0,1.78) {输出 $y$};
  \node[terminal] (finish) at (0,0.48) {结束};

  \draw[flowarrow] (start.south) -- (input.north);
  \draw (input.south) -- (loopjoin);
  \draw[flowarrow] (loopjoin) -- (test.north);
  \draw[flowarrow] (test.south) -- node[branchlabel,right=2pt] {是} (absupdate.north);
  % Affirmative branch loops left into the shared stem above the decision.
  \coordinate (loop-bottom) at (0,4.05);
  \coordinate (loop-left) at (-2.45,4.05);
  \coordinate (loop-left-top) at (-2.45,7.12);
  \coordinate (loop-arrow-tip) at (-0.35,7.12);
  \draw (absupdate.south) -- (loop-bottom) -- (loop-left) -- (loop-left-top);
  % The directed return path ends at the shared stem entry.
  \draw[flowarrow] (loop-left-top) -- (loop-arrow-tip) -- (loopjoin);

  % Negative branch bypasses the loop and enters the power update.
  \coordinate (right-branch) at (3.15,6.28);
  \coordinate (right-bottom) at (3.15,3.90);
  \coordinate (power-entry) at (0,3.90);
  \draw (test.east) -- node[branchlabel,above=2pt] {否} (right-branch) -- (right-bottom);
  \draw (right-bottom) -- (power-entry);
  \draw[flowarrow] (power-entry) -- (power.north);
  \draw[flowarrow] (power.south) -- (output.north);
  \draw[flowarrow] (output.south) -- (finish.north);

  \path[use as bounding box] (-2.85,0.15) rectangle (4.20,9.65);
\end{tikzpicture}
